\section{Contributions}

Elixir is currently taking a significant step towards integrating a type checker into its compiler.
Although the transition is gradual, both in its adoption and in the type system itself, the ultimate goal is to establish a static type system in Elixir.
To support this process, we make the following contributions in this paper:

\begin{itemize}
    \item \textit{Formalisation of TypeSpecs and Elixir Types.} 
    We formalise the existing TypeSpecs based on the TypeSpecs reference and the behaviour of Dialyzer. 
    The formalisation of Elixir Types is drawn from Elixir Type System by Giuseppe Castagna and Guillaume Duboc~\cite{duboc2024}.
    \item \textit{Translation from TypeSpecs to Elixir Types.} 
    We suggest the guideline to safely interpret Formalised TypeSpecs to Elixir Types.
    % \item Identification of mismatches between the two type systems.
    \item \textit{Approximating translation of undefined types.} 
    We identify types that are unable to directly translate and introduce conservative techniques to approximate the types.
    \item \textit{Prototype implementation of the proposed translation.}
    We apply the developed theory to real world data (public open-source repositories\footnote{The \num{175} source repositories, with the exact commits used during construction, are listed in the appendix (Table~\ref{tab:source-projects}).}) to produce a dataset that contains function definitions, original TypeSpecs, and translated types.
    \footnote{The translation prototype and dataset-construction pipeline are available at \url{https://github.com/KyleJYKim/type_migrator}.}
    \item \textit{Agreement study: Dialyzer versus the Elixir typechecker.}
    We measure the produced dataset by means of examining the validity and the quality of translation.
    \item \textit{LLM type prediction with translated Elixir Types.}
    We fine-tune a model (Qwen2.5+ 70b) based on the paper by H.B.Mengesha~\cite{Nina} with a filtered dataset that is the entries passed by both Dialyzer and the compiler.
    \footnote{The fine-tuning pipeline is available at \url{https://github.com/KyleJYKim/type_migrator_with_llm}.}
    The evaluation adopts a set-theoretic method by using \texttt{compatable?/2} from \texttt{Modulue.Types.Descr} developed by Duboc Guillaume.
    % on the checker-verified dataset (one on the precise, dynamic-free labels and one on the full both-pass set) to infer Elixir Types on untyped code, and compare them to assess whether admitting gradual types helps or hurts safe, precise prediction.
    % The intrinsic evaluation metric is based on semantic distance, a work done by H.B.Mengesha~\cite{Nina}, complemented by an extrinsic typecheck-acceptance metric.
\end{itemize}
